% ══════════════════════════════════════════════════════════════════
% LITERATURE REVIEW SECTION
% Insert after \section{Introduction} in balling_paper.tex
% ══════════════════════════════════════════════════════════════════

\section{Related Work}

\subsection{Balling Defect Formation Mechanisms}

Balling in LPBF arises from the competition between surface
tension and wetting forces along the melt track. The
Plateau--Rayleigh instability — the tendency of a liquid
cylinder to break into droplets to minimize surface energy —
is the dominant physical mechanism when the melt pool is
elongated relative to its diameter and solidification is slow
relative to the instability growth rate.
\citet{li2021phase} demonstrated through phase-field modeling
that fast solidification and substrate wettability both
counteract the balling effect, while high scan speed increases
the aspect ratio of the melt pool and promotes instability.
\citet{zhang2021understanding} applied Rayleigh--Plateau theory
directly to LPBF single tracks, showing that normalized enthalpy
determines whether a track is initially unstable and that dwell
time governs whether instability growth reaches detectable
amplitude. \citet{zoller2023numerical} used smoothed-particle
hydrodynamics to show that Marangoni and capillary forces drive
melt accumulation in the track tail region, which is the
initiating event for bead formation; similar ball size but
different morphology was found at different power settings.
A simple scaling model by \citet{lindstrom2023simple} showed
that the balling threshold can be expressed as a dimensionless
number combining material properties, powder size, and substrate
preheating, providing a rapid design tool for parameter
selection. Most recently, \citet{taylor2025analytical} developed
an analytical rivulet instability model accounting for substrate
stabilization, and found that solidification and instability
growth rates can be comparable — meaning balling is not a
purely fluid instability but is shaped by the coupled
thermal-fluid dynamics of the laser process. This body of
work establishes the physical basis for our observation that
balling wavelength is dominated by capillary effects rather than
power-dependent thermal factors.

\subsection{In-Situ Melt Pool Monitoring in LPBF}

In-situ monitoring of the melt pool using optical cameras and
photodiodes has become the principal approach for real-time
quality control in LPBF. Early work by
\citet{scime2018using} demonstrated that visible-light high-speed
cameras can capture melt pool morphology signatures indicative
of balling and keyholing in Inconel 718, and that unsupervised
learning on scale-invariant morphology descriptors can
differentiate between pool states without explicit labels.
The same group later showed that multi-scale CNNs can
autonomously detect powder spreading anomalies from layerwise
images \citep{scime2018multiscale}.

Camera-based approaches have since been extended to thermal
imaging. \citet{baumgartl2020deep} used thermographic off-axis
imaging combined with deep learning to detect delamination and
spatter with 96.8\% accuracy, demonstrating that non-contact
thermal sensing can complement visible-light cameras.
\citet{snow2021toward} trained CNNs on layerwise imagery with
X-ray CT as ground truth, achieving 93.5\% accuracy on large
flaws ($> 200\,\mu$m) and demonstrating cross-build
generalizability — an important practical requirement.
\citet{cao2023monitoring} used photodiode signals converted
to 2D images to classify single-track melting states into
three categories with 95.8\% accuracy using CNN.

For melt pool signal interpretation, \citet{reijonen2024effect}
showed that shielding gas flow speed has a highly non-linear
and direction-dependent effect on coaxial photodiode signals,
demonstrating that the gas flow geometry directly influences
what the sensor detects — a finding consistent with our
discovery of the Left/Right detectability asymmetry from
fixed-direction gas flow.

\subsection{Gas Flow, Spatter, and Melt Pool Asymmetry}

The interaction between shielding gas flow and the melt pool
is a critical but underappreciated source of process asymmetry.
\citet{esmaeilizadeh2018spatter} showed that spatter distribution
on the powder bed is governed by the interaction of melt pool
ejecta with the shielding gas flow field, and that surface
roughness increases significantly in spatter-rich regions of
the build plate — demonstrating spatial asymmetry driven by
fixed gas flow direction. \citet{zhang2020simulation} developed
a CFD model showing that the Coanda effect (gas flow tendency
toward the substrate) significantly affects spatter removal
patterns, with design-dependent spatial distributions.
\citet{li2022review} reviewed spatter generation mechanisms
and confirmed that spatter trajectory is governed jointly by
recoil pressure, surface tension, and shielding gas direction
— all of which are fixed relative to the scan direction.

The melt pool dynamics are further asymmetrized by the scan
direction. \citet{yuan2023revealing} showed through
DEM-CFD simulations that powder spatter follows three distinct
patterns (inward, upward, outward) depending on the relative
orientation of the vapor flow and shielding gas, with the
inward pattern specifically caused by shielding gas. This
provides a simulation-level basis for the asymmetry observed
in our experiments: the right (downwind) side of the track
receives cross-contamination from ejecta carried by the gas
flow, creating a more disturbed pool environment that produces
stronger and more consistent camera signals.

\subsection{Multi-Modal Sensor Fusion for Process Monitoring}

Single-modality monitoring has well-documented limitations —
no single sensor captures all relevant process phenomena.
\citet{petrich2021multimodal} presented one of the most
comprehensive multi-modal LPBF monitoring systems, fusing
layerwise imagery, acoustic emission, multi-spectral emission,
and scan vector data. Using neural networks on voxel-by-voxel
predictions, they achieved 98.5\% accuracy for binary flaw
classification with four-fold cross-validation, and showed
through sensitivity analysis that while optical imagery carries
the most information, additional modalities significantly improve
classification performance — directly motivating multi-modal
fusion.

\citet{gaikwad2022multi} fused coaxial high-speed video with
a temperature field imaging system, extracting physically
interpretable melt pool signatures (temperature, shape, spatter
intensity). The data fusion approach reduced false positive
rates from $\sim 0.1$ to $\sim 0.001$ without sacrificing
true positive rate, matching black-box deep learning approaches
but with interpretable features.
\citet{li2022convolutional} fused layerwise images, acoustic
emission, and photodiode signals through CNN at data, feature,
and decision levels, finding that feature-level fusion
outperformed all single-sensor baselines for SLM quality
monitoring. \citet{wu2024insitu} extended this to triple-sensor
(camera, photodiode, microphone) fusion, achieving 97.98\%
accuracy on high-quality samples.

A broader review by \citet{chen2024insitu} identified multi-sensor
fusion as the primary future direction for LPBF monitoring,
noting the need for standardization, improved reliability,
and decision-making strategies beyond early stopping. Our work
contributes to this direction by introducing a new fusion
modality — post-build surface profilometry — that has not
previously been combined with in-situ camera monitoring for
balling detection. Crucially, we show that this fusion is
non-trivial: the two modalities are nearly uncorrelated
($r \approx 0$), confirming they measure genuinely different
physical phenomena.

In-situ profilometry during the build process has also been
studied. \citet{spurek2024effect} used in-situ laser
profilometry to measure powder layer topography and correlate
it with part density, while \citet{zhang2023machine} combined
fringe projection profilometry with machine learning to measure
layer surface topography at improved accuracy and resolution.
Our work differs in using post-build profilometry as a
complementary sensing modality fused with camera data at the
frame level, rather than using in-situ profilometry as a
standalone quality metric.

\subsection{Deep Learning for Sequential Process Data}

Temporal modeling of melt pool sequences requires architectures
that can capture both spatial features within each frame and
temporal dynamics across frames. CNNs excel at spatial feature
extraction but treat each frame independently;
LSTM networks capture sequential dependencies but require
pre-extracted features as input. CNN-LSTM hybrid architectures
combine both capabilities and have shown strong results in
manufacturing anomaly detection. \citet{wang2020detecting}
demonstrated that RNN-based models can detect anomalies in
manufacturing time series by learning reconstruction errors
from normal operation data, a paradigm applicable to melt pool
monitoring. \citet{zhang2023machine} confirmed the utility of
hybrid CNN-LSTM architectures for layer-level quality monitoring
in polymer LPBF, achieving up to 99.7\% validation accuracy with
lightweight models.

In our work, the CNN-LSTM processes difference image sequences
rather than raw frames, which reduces the model's task to
detecting \emph{changes} in pool dynamics rather than absolute
pool state. This is motivated by the observation that balling
events manifest as transient perturbations — a brief disruption
of steady-state pool behavior — which are more visible in
frame differences than in raw images, particularly for lateral
bead types where the perturbation is asymmetric and transient.

\subsection{Vision-Language Models for Industrial Inspection}

Large vision-language models (VLMs) have demonstrated strong
zero-shot and few-shot generalization across visual tasks,
raising the question of whether domain-specific training is
necessary for defect detection. Recent work in industrial
inspection has explored this question: \citet{jeong2023winclip}
proposed WinCLIP, adapting CLIP for zero-shot anomaly detection
through compositional text prompts and multi-scale patch
features, achieving 91.8\% AUROC on MVTec-AD without any
training. \citet{li2024fade} similarly adapted CLIP for
few-shot industrial anomaly detection, significantly
outperforming task-specific trained models in zero-shot
settings. However, these methods target appearance-based
industrial defects (surface scratches, texture anomalies)
which differ substantially from the dynamic, temporal nature
of melt pool balling.

For domain-specific applications with highly specialized
visual signatures, the performance gap between VLMs and
task-trained models remains large. \citet{zhang2025seeing}
demonstrated this in wind turbine blade inspection, where
retrieval-augmented VLMs outperformed plain VLMs but the
improvement depended critically on the quality of the
retrieved domain context. The key challenge is that VLMs
trained on natural image-text data lack the domain knowledge
needed to interpret melt pool dynamics, spatter patterns,
and the temporal evolution of pool perturbations.

Our VLM evaluation confirms this picture for LPBF balling
detection: zero-shot performance is competitive only on
high-signal layers where the balling event is visually
obvious, and adding few-shot examples from other process
conditions \emph{hurts} performance because melt pool
appearance is strongly process-condition-dependent. This
directly parallels the finding of \citet{zhang2025seeing}
that domain context must match the test condition to be
helpful. The practical implication is that while VLMs may
serve as a rapid assessment tool or provide interpretable
outputs, task-trained models with explicit process parameter
conditioning remain necessary for reliable balling detection
across diverse process conditions. Future work with larger
datasets of labeled melt pool sequences — potentially
including fine-tuned VLMs or in-context learning with
process-matched examples — could narrow this gap.

\subsection{Positioning of This Work}

This paper makes several contributions not addressed by
prior work. First, existing melt pool monitoring approaches
for balling detection \citep{scime2018using, cao2023monitoring}
do not systematically analyze which bead categories are
detectable and why; we provide this analysis and attribute
the Left/Right detectability asymmetry to a specific physical
mechanism (gas flow geometry) rather than treating it as
a modeling artifact. Second, no prior work has fused
post-build profilometry with in-situ camera data for
balling detection; our multi-modal model is the first to
demonstrate that these two modalities are complementary
($r \approx 0$) and that fusion substantially improves
detection of previously undetectable bead categories.
Third, while multi-modal sensor fusion for LPBF monitoring
has been studied \citep{petrich2021multimodal,gaikwad2022multi},
prior work has fused co-temporal sensors (acoustic, optical,
thermal all recorded during the build); our fusion of
in-situ and post-build sensing introduces a new paradigm
where ex-situ measurements guide in-situ detection.
Finally, our process parameter ablation reveals that scan
speed $V$ conditioning can hurt generalization through
shortcut learning, a finding relevant to any model that
conditions on binary or low-cardinality process parameters.
\end{document}
